% !TEX program = xelatex

\documentclass[bachelor,UTF8,hyperref,oneside]{Template/scuthesis}

\usepackage{fancyvrb}
\usepackage{xcolor}
\usepackage{dirtree}
\usepackage{textcomp}
\renewcommand*{\DTstylecomment}{\color{blue}\kaishu}
\usepackage{supertabular}
\usepackage{makecell}
\usepackage{ltxtable}
\usepackage{multirow}
\usepackage{array}
\usepackage{booktabs}
\usepackage[draft]{changes}
\usepackage[AutoFakeBold]{xeCJK}
\usepackage{url}
\def\UrlBreaks{\do/\do-}

\usepackage[ruled,vlined,longend]{algorithm2e}
\SetAlgoLongEnd
\usepackage{listings}
\usepackage{fontspec}
\newfontfamily\consolas{Consolas}
\definecolor{codegreen}{RGB}{0,128,0}
\definecolor{codegray}{RGB}{120,120,120}
\definecolor{codepurple}{RGB}{160,32,240}
\definecolor{codeorange}{RGB}{200,120,0}
\definecolor{backcolor}{RGB}{245,245,245}
\lstset{
    language=Python,
    basicstyle=\consolas\fontsize{9pt}{11pt}\selectfont\color{black},
    keywordstyle=\color{codegreen}\bfseries,
    stringstyle=\color{codeorange},
    commentstyle=\color{codegray}\itshape,
    numberstyle=\color{codegray}\tiny,
    backgroundcolor=\color{backcolor},
    rulecolor=\color{codegray},
    breaklines=true,
    frame=single,
    showstringspaces=false,
    aboveskip=4pt,
    belowskip=4pt,
    lineskip=-1pt
}
\usepackage{tikz}
\usetikzlibrary{positioning, fit}

\usepackage{algpseudocode}
\usepackage{float}
\usepackage{amsmath}

\begin{document}
% 设置文档信息
\title{【页眉和摘要上的毕设论文题目】}
\ENGtitle{[English thesis title]}
\CoverTitle{【封面上的毕设论文题目】}

\author{【作者姓名】}
\Stunumber{【作者学号】}
\ENGauthor{[English name of the author]}
\accomplishdate{【完成时间】}
\collage{【学院名称】}
\grade{【20xx级】}
\supervisor{【指导教师姓名】}
\ENGsupervisor{[English name of the supervisor]}
\major{【专业名称】}
\ENGmajor{[English name of your major]}

\keywords{【中文关键词】}
\ENGkeywords{[English keywords]}
% 设置论文正文前的页码、页眉等
\frontmatter\pagestyle{fancy}\makechaptertitlecenter
% 自动制作标题
\maketitle
%!TEX root = ../MainBody.tex

\begin{CHSabstract}
    \chabstract{【中文摘要。】}
\end{CHSabstract}

\begin{ENGabstract}
    \enabstract{[English abstract.]}
\end{ENGabstract}

% 自动制作目录
\maketoc
% 设置论文正文部分的页码、页眉等
\mainmatter\pagenumbering{arabic}\pagestyle{fancy}\makechaptertitlecenter
% 添加论文的各个章节，因为最终论文可能很长，把所有章节都写在一个文件里面会很长，用这种方法可以缩短篇幅，看起来更清楚
%!TEX root = ../MainBody.tex

\chapter{绪论}

\section{二级标题}

\subsection{三级标题}

%!TEX root = ../MainBody.tex

\chapter{相关研究情况和本文研究概述}

引用示例：ResNet~\cite{he2016deep}

\section{二级标题}

\subsection{三级标题}

%!TEX root = ../MainBody.tex

\chapter{【毕设主要内容名称】}

\section{二级标题}

\subsection{三级标题}

%!TEX root = ../MainBody.tex

\chapter{实验}

\section{二级标题}

\subsection{三级标题}

%!TEX root = ../MainBody.tex

\chapter{实验结果和分析}

\section{二级标题}

\subsection{三级标题}

%!TEX root = ../MainBody.tex

\chapter{【系统实现（如有）】}

\section{二级标题}

\subsection{三级标题}

%!TEX root = ../MainBody.tex

\chapter{总结和展望}

\section{二级标题}

\subsection{三级标题}

% 设置论文正文后的式样
\backmatter\makechaptertitlecenter
% 按国标自动制作参考文献
\begin{reference}
    \bibliographystyle{Template/chinesebst}\bibliography{Reference}
\end{reference}
% 致谢
%!TEX root = ../MainBody.tex

\chapter{致谢}

% 声明与版权使用授权书
\include{Chapters/Combined_Authorization}
% AI工具使用声明
%!TEX root = ../MainBody.tex

\chapter{四川大学本科毕业论文 AI 工具使用声明}

\vspace{1em}

\noindent\hspace{2em}本人声明，在撰写题目为《【毕设论文题目】》的本科毕业论文过程中，对 AI 工具的使用情况如下：

\vspace{0.8em}

\noindent\textbf{一、使用的 AI 工具及使用目的}

\vspace{0.5em}

\begingroup
\renewcommand{\arraystretch}{1.35}
\setlength{\tabcolsep}{5pt}
\begin{center}
    \begin{tabular}{>{\centering\arraybackslash}p{1cm} >{\centering\arraybackslash}p{5cm} >{\centering\arraybackslash}p{8cm}}
        \toprule
        \centering 序号 & \centering AI 工具 & \centering\arraybackslash 使用目的 \\
        \midrule
        \centering 1 & 【AI工具名称】 & 【用AI工具做了什么】 \\
        \midrule
        \centering \ldots & \ldots & \ldots \\
        \bottomrule
    \end{tabular}
\end{center}
\endgroup

\vspace{0.8em}

\noindent\textbf{二、使用程度说明}

\vspace{0.5em}

\begingroup
\renewcommand{\arraystretch}{1.35}
\setlength{\tabcolsep}{4pt}
\begin{center}
    \begin{tabular}{>{\centering\arraybackslash}p{1cm} >{\centering\arraybackslash}p{2cm} >{\centering\arraybackslash}p{3cm} >{\centering\arraybackslash}p{3cm} >{\centering\arraybackslash}p{2cm} >{\centering\arraybackslash}p{3cm}}
        \toprule
        \centering 序号 & \centering 章节 & \centering AI 工具 & \centering 使用方式 & \centering 内容占比 & \centering\arraybackslash 处理方式 \\
        \midrule
        \centering 1 & 第一章 & 【AI工具】 & 【用AI工具做了什么】 & 【百分比】 & 【怎么用的】 \\
        \midrule
        \centering \ldots & \ldots & \ldots & \ldots & \ldots & \ldots \\
        \bottomrule
    \end{tabular}
\end{center}
\endgroup

\vspace{1em}

\noindent\textbf{三、原创性与责任声明}

\vspace{0.5em}

\noindent\hspace{2em}本人确认，尽管使用了上述 AI 工具，但论文整体的研究思路、核心观点、论证过程及最终结论均为本人独立思考与研究的成果。论文中引用 AI 工具生成内容之处，均已按照学校规定的学术规范进行了明确标注。

\vspace{0.6em}

\noindent\hspace{2em}若论文因 AI 工具使用不当或未按要求标注等原因出现学术不端问题，本人愿意承担全部责任。

\vspace{3em}

\begin{flushright}
    \begin{tabular}{@{}l@{}}
        学生签名：\hspace{4cm}\\[1.2em]
        指导老师签名：\hspace{4cm}\\[1.2em]
        【日期】 \hspace{4cm}
    \end{tabular}
\end{flushright}
\end{document}
